\documentclass{style/presentation}

\title{Identification et optimisation énergétique pour l'usinage}
\author{Jamy \textsc{Vignaud}}
\institute{IMS, I2M, VLM Robotics}
\date{2 décembre 2026}

\begin{document}

\frame{\titlepage}

\section{Introduction}

\subsection{Contexte industriel et enjeux}

\begin{frame}
    VLM Robotics, contraintes résiduelles, minimisation de l'énergie
\end{frame}

\subsection{Problématique et contributions}

\begin{frame}
    Verrou de la boucle fermée (identification) + Minimisation de l'énergie (optimisation)

    Plan
\end{frame}

\begin{frame}[sanssoustitre]{Plan}
    \tableofcontents
\end{frame}

\section{Moyens expérimentaux}

\subsection{Machine-outil et chaîne d'acquisition}

\begin{frame}

\end{frame}

\subsection{Étalonnage de la platine dynamométrique par un cobot}

\begin{frame}

\end{frame}

\section{Modélisation mécanique de l'énergie de coupe}

\subsection{Modèle de puissance par torseurs}

\begin{frame}
    Comoment, cinématique, actions de la dent
\end{frame}

\subsection{Identification des pressions spécifiques}

\begin{frame}
    Traitement du signal, Kienzle-Victor
\end{frame}

\subsection{Comparaison à la puissance électrique}

\begin{frame}

\end{frame}

\subsection{Limite}

\begin{frame}
    $W_c$ indépendant de $V_c$
\end{frame}

\section{Identification en boucle fermée de la broche}

\subsection{Problème et positionnement}

\begin{frame}
    architecture Sinamics, état de l'art
\end{frame}

\subsection{Algorithmes proposés}

\begin{frame}
    LS -> IV -> SRIV -> RIV
\end{frame}

\subsection{Validation en simulation}

\begin{frame}
    Monte-Carlo, bruit blanc puis coloré
\end{frame}

\subsection{Validation expérimentale}

\begin{frame}
    filtre stabilisateur, résultat sur machine, points de fonctionnement
\end{frame}

\section{Optimisation énergétique en ligne}

\subsection{Modèle empirique et algorithme RPEM}

\begin{frame}

\end{frame}

\subsection{Validation en simulation}

\begin{frame}

\end{frame}

\subsection{Résultats sur machine-outil}

\begin{frame}

\end{frame}

\section{Conclusion et perspectives}

\begin{frame}

\end{frame}

\end{document}